% =============================================================================
% macros.tex: shared by every main-*.tex; \input it right after \documentclass
% (and the class's own front-matter set-up).
%
% \TODOOPTS decides whether todonotes show. The Makefile sets it:
%   make pdf    \def\TODOOPTS{disable}  clean PDF, every \todo hidden
%   make todos  \def\TODOOPTS{}         margin notes and \listoftodos
% Compiled any other way (Overleaf, an editor), the todos show.
% =============================================================================
\providecommand{\TODOOPTS}{}
\usepackage[\TODOOPTS,colorinlistoftodos,textsize=footnotesize]{todonotes}
\setlength{\marginparwidth}{2cm}
\usepackage{amsmath,graphicx,booktabs,catchfile,xcolor}
\ifdefined\Bbbk\else\usepackage{amssymb}\fi % asmejour's newtxmath already has the AMS symbols
% max width: a 175 mm figure shrinks to fit a narrower review layout; at the
% journal's final layout it is placed at its own size.
\usepackage[export]{adjustbox}
\usepackage{siunitx}
\sisetup{separate-uncertainty = true}

% Figures come only from figures/out/, written by make figures.
\graphicspath{{../figures/out/}}
\newcommand{\vclfigdir}{../figures/out/}

% \figcaption{fig1_slug}: the caption stub figures/fig1_slug.py wrote, which
% states n and what the error bars are. Edit the words in the script.
\newcommand{\figcaption}[1]{%
  \CatchFileDef{\vclcaption}{\vclfigdir#1.caption.tex}{}\caption{\vclcaption}}

% \fillme{...}: something to fill in from the record, never to invent. Violet
% in BOTH builds, so a PDF that leaves the lab cannot hide one;
% make submission-check lists every one left.
\definecolor{fillme}{RGB}{142,68,173}
\NewDocumentCommand{\fillme}{m}{\textcolor{fillme}{[#1]}}
% \dummy{1.23}: a stand-in number, violet in both builds (tensegrity PR 76).
\NewDocumentCommand{\dummy}{m}{\textcolor{fillme}{#1}}

% Placeholders for a figure or table that has no script or data yet; both
% vanish into a todo in the review build (from tensegrity PR 76).
\newcommand{\figplaceholder}[2]{% #1 label, #2 what it will show
  \begin{figure}[ht]
    \centering
    \fbox{\parbox[c][3.2cm][c]{0.85\linewidth}{\centering\textit{Figure placeholder: #2}}}%
    \caption{\fillme{Placeholder.} #2}
    \label{fig:#1}
  \end{figure}%
  \todo[inline,color=orange!30]{Figure \texttt{#1}: #2}}
\newcommand{\tabplaceholder}[2]{% #1 label, #2 what it will hold
  \begin{table}[ht]
    \centering
    \caption{\fillme{Placeholder.} #2}
    \label{tab:#1}
    \begin{tabular}{@{}lll@{}}
      \toprule
      Column A & Column B & Column C \\
      \midrule
      \multicolumn{3}{c}{\textit{(table contents pending)}} \\
      \bottomrule
    \end{tabular}
  \end{table}%
  \todo[inline,color=orange!30]{Table \texttt{#1}: #2}}
